% !TeX root = ../main.tex
% ============================================
% 第二章 相关理论与技术基础
% ============================================

\chapter{相关理论与技术基础}

\section{多模态数据融合理论}

\subsection{非完全观测下的健康状态估计（Observer）}
自动化健康管理系统的控制对象是个体健康状态 $x_t$。该状态无法被直接完整测量，只能通过多源传感器与业务接口形成观测量 $y_t$。因此，本研究将健康评估建模为非完全观测系统（Partial Observability），并采用状态方程与观测方程描述闭环过程：
\begin{equation}
    x_{t+1}=f(x_t,u_t,w_t),
    \label{eq:state_transition_21}
\end{equation}
\begin{equation}
    y_t=h(x_t)+v_t,
    \label{eq:measurement_model_21}
\end{equation}
其中，$u_t$ 为干预输入，$w_t$ 与 $v_t$ 分别表示过程噪声和观测噪声。该建模解释了采用贝叶斯推断的必要性：系统目标是依据不完整观测估计潜在状态及其风险分布，而非对原始观测做直接映射。

系统当前观测通道由以下四类组成：
\begin{itemize}
\item 结构化画像观测：\texttt{/user/profile}。
\item 文档观测：\texttt{/api/v1/ocr/upload} 将体检单解析为结构化字段。
\item 流式生理观测：\texttt{/api/v1/iot/sync/batch} 提供高频心率等实时信号。
\item 决策入口观测融合：\texttt{/analyze/comprehensive} 对数据库画像与请求增量做合并。
\end{itemize}

\subsection{证据强度与时间尺度分层的贝叶斯融合}
在风险决策层，临床先验 $P_c$、基因修正系数 $\alpha_t$、行为修正系数 $\beta_t$ 按乘性方式完成融合：
\begin{equation}
    P_t=\operatorname{clip}(P_c \cdot \alpha_t \cdot \beta_t,\;0.1,\;99.9).
    \label{eq:fusion_21}
\end{equation}
在后续章节中，当前时刻的融合风险也记作 $P_{fusion}$；对于单次静态决策时刻，二者含义等价。进一步地，$\alpha_t$ 与 $\beta_t$ 被定义为证据调节因子（Evidence Modifiers）。代码实现中，$\alpha_t\in[0.8,1.5]$，$\beta_t\in[0.8,1.3]$，分别对应遗传证据与行为证据对临床先验的调节强度。

为刻画证据贡献的可解释性，引入未截断风险 $\tilde{P}_t=P_c\alpha_t\beta_t$，可写为：
\begin{equation}
    \log \tilde{P}_t=\log P_c+\log \alpha_t+\log \beta_t.
    \label{eq:evidence_strength_21}
\end{equation}
该表达表明，融合贡献在对数域可加，便于比较不同证据通道的边际作用。时间尺度上，基因证据属于慢变量，行为证据属于快变量，两者与临床先验形成多时间尺度层级耦合。

\subsection{安全约束与后验更新机制}
系统在静态融合后引入实时后验修正，以事件因子 $\gamma_t$ 表达短时扰动：
\begin{equation}
    P_{t+1}=\operatorname{clip}(\gamma_t\cdot P_t,\;0.1,\;99.9).
    \label{eq:online_update_21}
\end{equation}
实时更新规则可形式化为：
\begin{equation}
\gamma_t=
\begin{cases}
1.8, & HR_t>100 \land BMI>28, \\
1.3, & HR_t>100 \land BMI\leq 28, \\
1.2, & HR_t<50, \\
1.0, & \text{otherwise}.
\end{cases}
\label{eq:gamma_rule_21}
\end{equation}
此外，为降低高遗传易感个体漏检风险，系统引入潜在风险拦截规则：
\begin{equation}
\text{Level}=\text{Potential},\quad
\text{if }\alpha_t\geq 1.2\;\land\;P_t\leq 20.
\label{eq:potential_rule_21}
\end{equation}
上述截断与拦截共同构成安全约束层，分别用于保障概率可解释边界与高危亚群识别敏感性。

\subsection{理论-工程映射与闭环一致性}
本节公式与工程实现保持一一对应关系。本节式~\eqref{eq:fusion_21} 与第四章核心融合式~\eqref{eq:fusion}同构，均对应 \texttt{fusion\_service.py} 的主融合路径。映射关系如表~\ref{tab:theory_code_mapping_21} 所示。

\begin{table}[htbp]
\centering
\caption{理论公式与工程实现映射关系}
\label{tab:theory_code_mapping_21}
\begin{tabular}{p{2.6cm}p{4.2cm}p{3.4cm}p{2.4cm}}
\toprule
理论表达 & 代码落点 & 接口链路 & 功能作用 \\
\midrule
式~\eqref{eq:fusion_21} & \texttt{fusion\_service.py}\\\texttt{calculate\_composite\_risk} & \texttt{/analyze/}\\\texttt{comprehensive} & 静态融合决策 \\
式~\eqref{eq:online_update_21} & \texttt{fusion\_service.py}\\\texttt{update\_realtime\_risk} & \texttt{/api/v1/iot/}\\\texttt{sync/batch} & 实时后验更新 \\
式~\eqref{eq:potential_rule_21} & \texttt{gene\_mod>=1.2}\\\texttt{\& risk<=20} 分支判定 & 融合引擎内部规则 & 漏检抑制 \\
式~\eqref{eq:measurement_model_21} & Profile/OCR/IoT 三通道观测汇聚 & \texttt{/user/profile}\\\texttt{/api/v1/ocr/upload} & 观测重建 \\
\bottomrule
\end{tabular}
\end{table}

从闭环控制角度看，系统形成了“观测（Profile/OCR/IoT）-决策（贝叶斯融合）-控制（干预建议）-反馈（实时更新）”的结构化链路，为第四章算法推导、第五章实现细节和第六章实验分析提供统一理论前提。

\section{机器学习与深度学习算法}

\subsection{临床风险预测模型（LightGBM）}
本系统将慢病风险预测建模为“多病种并行二分类”任务。训练脚本以每个 \texttt{Target\_*} 标签为独立学习目标，形成模型集群，并在推理阶段输出临床先验风险 $P_c$。

模型采用梯度提升树的加法框架：
\begin{equation}
    F_M(\mathbf{x})=\sum_{m=1}^{M} \eta h_m(\mathbf{x}),
    \label{eq:lgb_additive_22}
\end{equation}
\begin{equation}
    P_c=\sigma\!\left(F_M(\mathbf{x})\right)\times 100\%.
    \label{eq:clinical_prior_22}
\end{equation}
其中，$\eta$ 为学习率，$h_m$ 为第 $m$ 棵基学习树。该表达对应 \texttt{train\_risk\_models.py} 的训练与概率输出流程。

工程实现中，数据划分采用分层抽样
\texttt{train\_test\_split(..., test\_size=0.2, stratify=y)}，并启用
\texttt{is\_unbalance=True} 处理类别不均衡。为抑制标签泄漏，针对病种 $d$ 的训练特征集合定义为：
\begin{equation}
    \mathcal{X}_d = \mathcal{X}_{all} \setminus \mathcal{L}_d,
    \label{eq:leakage_control_22}
\end{equation}
其中 $\mathcal{L}_d$ 由 \texttt{LEAKAGE\_MAPPING} 动态给出。该机制保证“诊断同义特征”不直接进入模型输入。

\subsection{行为风险分类与证据调节因子（XGBoost）}
行为侧模型以 \texttt{sum}、\texttt{count} 两个统计特征预测 \texttt{Lifestyle\_Risk}，并以 Accuracy、Precision、Recall、F1 作为分类评估指标。该模型在系统中的角色不是独立终判，而是生成融合引擎的行为证据调节因子 $\beta$。

在线推理时，\texttt{lifestyle\_service.py} 先将步数映射到训练域：
\begin{equation}
    c_{sim}=C_0,\qquad s_{sim}=\frac{steps}{S_0}\,C_0,
    \label{eq:steps_mapping_22}
\end{equation}
其中 $S_0=10000$、$C_0=1000$。随后由分类器输出高风险概率 $p_{life}$，并映射为行为修正系数：
\begin{equation}
    \beta = b_0 + b_1 p_{life},\quad b_0=0.8,\; b_1=0.5,
    \label{eq:beta_mapping_22}
\end{equation}
即 $\beta\in[0.8,1.3]$。该定义与融合引擎的乘性证据结构保持一致。

\subsection{视觉多任务营养回归模型（EfficientNet）}
视觉模型采用 EfficientNet-B0 作为骨干网络，并将分类头替换为四维回归输出，直接预测
\{Calories, Carbs, Protein, Fat\}。该选型理由在于 Compound Scaling 同时调节网络深度、宽度和输入分辨率，在精度与计算开销之间取得可控折中：
\begin{equation}
    d=\alpha^{\phi},\qquad w=\beta^{\phi},\qquad r=\gamma^{\phi},
    \label{eq:compound_scaling_22}
\end{equation}
其中满足约束 $\alpha\beta^2\gamma^2\approx 2$。该性质使模型更适合实时交互场景下的营养估计任务。

训练目标采用多任务 L1 损失：
\begin{equation}
    \mathcal{L}_{\mathrm{L1}}=\frac{1}{4}\sum_{k=1}^{4}\left|\hat{y}_k-y_k\right|,
    \label{eq:vision_l1_22}
\end{equation}
并在服务端推理阶段执行非负约束
\begin{equation}
    \hat{y}_k^{+}=\max(0,\hat{y}_k),
    \label{eq:vision_nonneg_22}
\end{equation}
保证营养值物理可解释。模型路径由配置项 \texttt{NUTRITION\_MODEL\_PATH} 统一管理，默认指向 \texttt{nutrition\_efficientnet.pth}。

\subsection{基于反事实推理的干预决策支持（Counterfactual Reasoning）}
本系统的“平行宇宙推演”采用反事实推理范式。理论上，可将健康系统抽象为因果图
$\mathcal{G}=(\mathcal{V},\mathcal{E})$，通过对干预节点执行结构赋值，比较事实世界与干预世界的风险差异，而非仅做相关性外推。

在工程实现中，\texttt{projection\_service.py} 提供两类推演：
\begin{itemize}
\item 自然演化推演：按年递增年龄并对 SBP、空腹血糖、BMI 施加固定漂移。
\item 干预推演：对体重施加 \texttt{weight\_loss\_percent}，并联动更新 BMI、SBP、\texttt{Glucose\_Fasting}。
\end{itemize}
干预后的状态修改可表示为：
\begin{equation}
\mathbf{x}^{cf}=\mathrm{do}\!\left(
W\leftarrow(1-\lambda)W,\;
SBP\leftarrow SBP-5,\;
G\leftarrow G-0.5
\right),
\label{eq:cf_intervention_22}
\end{equation}
其中 $\lambda$ 为减重比例，$G$ 表示空腹血糖。事实世界与反事实世界均复用同一融合引擎计算风险，收益定义为：
\begin{equation}
    \Delta_{risk}=\frac{P^{base}_{max}-P^{new}_{max}}{P^{base}_{max}}\times 100\%.
    \label{eq:delta_risk_22}
\end{equation}

\subsection{算法与数据流映射（Data Flow Graph）}
为连接第 2 章理论与第 5/6 章工程验证，本节给出统一数据流主线：
\begin{equation}
    P_c \xrightarrow[\text{gene/lifestyle}]{\alpha,\beta}
    P_{fusion} \xrightarrow[\text{counterfactual}]{\mathcal{I}}
    \Delta_{risk}.
    \label{eq:dataflow_22}
\end{equation}
其中 $\mathcal{I}$ 表示干预算子集合。对应实现路径如下：
\begin{itemize}
\item \texttt{train\_risk\_models.py} 与 \texttt{risk\_engine.py} 生成并调用临床先验 $P_c$。
\item \texttt{gene\_service.py} 与 \texttt{lifestyle\_service.py} 分别提供 $\alpha$ 与 $\beta$。
\item \texttt{fusion\_service.py} 执行概率连乘与安全截断，输出 $P_{fusion}$。
\item \texttt{projection\_service.py} 在事实/反事实双路径上计算 $\Delta_{risk}$。
\end{itemize}
该映射使算法链路与控制闭环保持统一：感知输入经融合决策形成风险状态，再通过干预算子生成可量化收益。
\section{自然语言处理与知识增强}

\subsection{检索增强生成的知识索引与向量检索机制}
本系统采用“离线索引构建 + 在线语义检索”的两阶段知识增强流程。离线阶段将指南类 PDF 文档加载为页面文本，随后执行递归分块，分块参数为 \texttt{chunk\_size=500}、\texttt{chunk\_overlap=50}；在线阶段对用户查询执行 Top-$k$ 相似检索并返回证据片段。

知识分块可表示为：
\begin{equation}
    \mathcal{C}=\operatorname{Split}(\mathcal{D};L=500,o=50),
    \label{eq:rag_chunk_23}
\end{equation}
其中 $\mathcal{D}$ 为文档集合，$\mathcal{C}$ 为分块集合，$L$ 与 $o$ 分别为块长与重叠长度。向量检索定义为：
\begin{equation}
    \mathcal{R}(q)=\operatorname{TopK}_{c_i\in\mathcal{C}}\,\operatorname{sim}(e(q),e(c_i)),\quad K=3,
    \label{eq:rag_retrieve_23}
\end{equation}
其中 $e(\cdot)$ 为嵌入映射，\(\operatorname{sim}(\cdot)\) 为相似度函数。系统实现采用 \texttt{HuggingFaceEmbeddings(text2vec-base-chinese)} 与 \texttt{Chroma} 持久化向量库。

\subsection{用户画像约束下的受控生成机制}
对话模块并非纯生成流程，而是由“查询语义 + 检索证据 + 用户画像”共同约束的受控生成。其形式化表达为：
\begin{equation}
    a=\mathcal{G}\!\left(q,\mathcal{R}(q),p_u\right),
    \label{eq:chat_controlled_23}
\end{equation}
其中 $q$ 为用户问题，$\mathcal{R}(q)$ 为检索证据，$p_u$ 为用户画像，$a$ 为最终回答。

工程实现中，\texttt{chat\_service} 先生成 \texttt{profile\_summary}，再调用 \texttt{rag\_service.search\_context(query, k=3)}，并将两者拼接至提示词后提交 LLM 生成。该链路保证回答在医学知识约束下兼顾个体差异，同时返回 \texttt{sources} 作为证据追溯接口。

\subsection{多级置信度验证架构（Multi-stage Confidence-based Verification）}
医疗文档解析采用“OCR 感知—LLM 结构化—规则校验兜底”的多级验证架构，而非单一模型直出。其目的在于提升事实一致性并降低解析失真风险。

对输入医学文档 $x$，定义 OCR 输出为 $z_{ocr}=f_{ocr}(x)$，结构化结果为：
\begin{equation}
\hat{z}=
\begin{cases}
    f_{llm}(z_{ocr}), & s_{llm}=1,\\
    f_{regex}(z_{ocr}), & s_{llm}=0,
\end{cases}
\label{eq:multi_stage_verify_23}
\end{equation}
其中 $s_{llm}$ 表示 LLM 路径可用性（含 API 成功、JSON 可解析等判据）。

系统实现包含三项关键鲁棒机制：
\begin{itemize}
\item PDF 预处理：单次任务最多处理前 5 页（\texttt{max\_pages=5}），并将页面渲染为高分辨率图像。
\item 并发 OCR：采用 \texttt{ThreadPoolExecutor(max\_workers=5)} 并发识别多页文本。
\item 重试与回退：LLM 路径使用 \texttt{tenacity} 指数退避重试（最多 3 次）；失败后自动切换正则抽取。
\end{itemize}
重试等待时间可抽象为：
\begin{equation}
    \Delta t_n=\min\left(10,\max\left(2,2^{n-1}\right)\right),\quad n=1,2,3.
    \label{eq:retry_wait_23}
\end{equation}

\subsection{事件驱动的数据一致性保障机制}
本系统将缓存层定义为一致性控制组件，而非仅用于性能优化。缓存键采用用户粒度命名规范，核心模式包括
\texttt{chat\_response:\{user\_id\}:*}、\texttt{diet\_plan:\{user\_id\}:*} 与
\texttt{risk\_assessment:\{user\_id\}:*}。

当发生关键状态事件（如 \texttt{/user/profile} 更新或 \texttt{/api/v1/ocr/upload} 导入新报告）时，系统触发 \texttt{invalidate\_user\_cache(user\_id)} 执行模式删除。其一致性规则可表示为：
\begin{equation}
    E_t(u)=1\Rightarrow \mathcal{K}_{t+1}(u)=\varnothing,
    \label{eq:event_consistency_23}
\end{equation}
其中 $E_t(u)$ 为用户 $u$ 的状态变更事件，$\mathcal{K}_{t+1}(u)$ 为变更后保留缓存集合。对话链路额外提供 \texttt{force\_refresh} 旁路控制，以显式跳过缓存并回源生成。

\subsection{感知层质量边界与闭环控制映射}
在自动化闭环中，NLP 与 OCR 模块承担“状态变量生成器”的角色。其输出质量直接影响决策层风险估计与控制层干预建议。

设真实状态为 $x_t$，感知估计为 $\hat{x}_t$，则感知误差
$\varepsilon_{obs}=\|\hat{x}_t-x_t\|$ 与决策误差 $\varepsilon_{dec}$ 存在单调耦合关系：
\begin{equation}
    \varepsilon_{dec}=g(\varepsilon_{obs}),\quad g'(\cdot)>0.
    \label{eq:gigo_23}
\end{equation}
该关系给出本节的系统学意义：通过 RAG 证据约束、OCR 多级验证与事件驱动一致性控制，感知层对下游决策层形成“误差抑制闸门”，以降低 GIGO（Garbage In, Garbage Out）效应。
\section{系统关键技术栈}

\subsection{服务化架构、控制反转与可测试性}
本系统的工程实现不以单体脚本直接堆叠算法模块，而是采用“接口适配层--业务服务层--数据持久层”的服务化分层。路由层负责 HTTP 协议适配、鉴权校验与输入输出规范化；服务层封装风险评估、OCR 解析、RAG 问答与干预推演等业务算子；数据层通过统一会话接口完成状态读写。该结构使感知、决策与反馈模块能够在稳定接口约束下协同运行。

控制反转的核心在于：业务函数不直接创建其依赖，而是在运行时接收数据库会话、用户上下文或服务实例。若将业务函数记为 $f$，则可测试性可抽象为：
\begin{equation}
    y=f(x;d_1,\ldots,d_n),\qquad
    \hat{y}=f(x;\hat{d}_1,\ldots,\hat{d}_n),
    \label{eq:testability_24}
\end{equation}
其中 $d_i$ 为运行态依赖，$\hat{d}_i$ 为测试态替代依赖。当输入 $x$ 保持不变时，输出变化可主要归因于依赖替换而非业务逻辑本体，从而提高系统的可观测性与可验证性。

该理论在工程中体现为以下三类机制：
\begin{itemize}
\item 运行态注入：后端接口普遍通过 \texttt{Depends(get\_session)}、\texttt{Depends(get\_current\_user)} 等方式接收数据库会话与认证上下文，而非在函数内部直接实例化依赖。
\item 生命周期解耦：核心模型与缓存连接在 \texttt{lifespan} 阶段统一初始化，由此实现请求处理路径与底层资源加载过程的分离。
\item 测试态替换：测试环境采用 \texttt{dependency\_overrides} 覆盖真实会话提供器，并结合内存 SQLite 与 \texttt{StaticPool} 构造隔离数据库；对重型模型服务，则通过模块级 Mock 抑制无关基础设施对接口测试的干扰。
\end{itemize}
因此，控制反转在本系统中不仅是软件解耦手段，也是一种面向系统测试的工程组织原则。第六章中接口隔离测试与端到端场景验证，均以前述可测试性结构为实现前提。

\subsection{异步非阻塞I/O、排队稳定性与缓存一致性}
个性化健康管理平台同时面向 OCR 上传、RAG 问答、IoT 数据同步与风险更新等多类异构请求。此类请求的共同特征在于：计算负载与外部 I/O 等待交替出现，若采用同步阻塞式处理，请求队列长度将随并发上升而快速累积。为分析该问题，可引入 Little 定律：
\begin{equation}
    L=\lambda W,
    \label{eq:little_24}
\end{equation}
其中 $L$ 为系统内平均请求数，$\lambda$ 为平均到达率，$W$ 为平均逗留时间。在到达率既定时，保持系统稳态的关键在于抑制 $W$ 的增长。

对单次请求而言，其平均逗留时间可进一步分解为：
\begin{equation}
    W=t_{queue}+t_{io}+t_{compute},
    \label{eq:latency_24}
\end{equation}
其中 $t_{queue}$ 表示排队等待时间，$t_{io}$ 表示外部模型调用、文件读取与缓存访问等 I/O 等待时间，$t_{compute}$ 表示本地推理与规则计算时间。异步非阻塞 I/O 的作用并非消除服务需求，而是在 I/O 等待阶段释放执行上下文，从而降低等效排队长度并抑制系统拥塞。

本系统中与式（\ref{eq:little_24}）对应的实现主要包括：
\begin{itemize}
\item 对话与 OCR 服务均采用异步客户端调用外部大模型接口，减少长时网络等待对主请求链路的占用。
\item 知识库重建与部分管理任务通过 \texttt{BackgroundTasks} 后台调度执行，将高耗时作业从实时响应路径中分离。
\item 医疗文档解析在异步主流程之外引入 \texttt{ThreadPoolExecutor} 并发 OCR 与指数退避重试，以降低多页文档处理中的等待放大效应。
\end{itemize}

缓存层则承担一致性控制功能，而非单纯的性能加速角色。系统通过 \texttt{CacheManager} 提供异步读写、模式匹配删除与用户粒度失效机制；当用户档案更新或新体检报告导入后，系统立即触发 \texttt{invalidate\_user\_cache}，使缓存状态与用户健康状态保持同步。该机制与式（\ref{eq:event_consistency_23}）共同构成事件驱动的一致性控制链路；同时，对话接口保留 \texttt{force\_refresh} 旁路，以支持在高敏感场景下显式回源计算。

\subsection{基础设施即代码与实验可复现性}
对于涉及多模型推理、文件存储与缓存协同的系统而言，实验结果不仅受算法参数影响，也受运行环境差异影响。为抑制环境扰动，本系统采用容器化方式固定关键基础设施，并将服务拓扑、镜像版本、网络与卷映射纳入统一编排。实验环境可抽象为：
\begin{equation}
    \mathcal{E}=(\mathcal{I},\mathcal{S},\mathcal{V},\mathcal{N},\mathcal{C}),
    \label{eq:repro_24}
\end{equation}
其中 $\mathcal{I}$ 为镜像集合，$\mathcal{S}$ 为服务集合，$\mathcal{V}$ 为持久化卷映射，$\mathcal{N}$ 为网络拓扑，$\mathcal{C}$ 为关键配置项。通过固定 $\mathcal{E}$，可将实验误差中由基础设施变化引入的非目标扰动压缩至可控范围。

本系统的容器编排采用 \texttt{docker-compose.yml} 统一定义后端 API、前端 SPA 与 Redis 缓存三类服务；后端镜像基于 \texttt{python:3.11-slim}，前端采用 \texttt{node:20-alpine} 构建并以 \texttt{nginx:alpine} 提供静态资源服务。编排文件进一步固定了桥接网络、健康检查与卷挂载策略，并通过环境变量显式注入 \texttt{REDIS\_URL} 等关键运行参数；视觉模型路径则由配置项 \texttt{NUTRITION\_MODEL\_PATH} 统一管理。

需要指出的是，当前容器化方案的主要目标是实验环境复现与依赖封装，而非生产级弹性调度。应用启动入口仍保留开发态热重载设置，且数据库路径由应用配置统一管理。因此，本节讨论的基础设施即代码应理解为“最小可复现实验部署”，其作用在于为第五章实现验证与第六章实验复现提供稳定环境边界。


